% §6 Related Work

\section{Related Work}
\label{sec:related}

\paragraph{Vault custody protocols.}
Swambo et al.~\cite{SHMB20} formalised the vault lifecycle---deposit,
unvault, withdraw, recover---and the watchtower monitoring assumption
that underpins all covenant vault designs, establishing the threat
vocabulary we adopt and extend to the 11-element taxonomy. Their
companion work~\cite{SBMB20} characterised three implementation
approaches (deleted-key pre-signed transactions, COSHV, \opctv{}) but
did not compare opcode-level designs empirically. Swambo's doctoral
thesis~\cite{Swa23} extended this to the Ajolote autonomous custody
system but did not include \opctv{}, \opccv{}, \opvault{}, or
\catcsfs{}. Swambo and Poinsot's risk framework for the Revault
protocol~\cite{SP21} developed an operational risk taxonomy for
pre-signed-transaction vaults; our per-threat analysis is in a
similar spirit but targets opcode-level covenants rather than
deleted-key designs. Möser, Eyal, and Gün Sirer~\cite{MES16}
introduced covenants as a vault-custody mechanism, predating all
four BIPs. O'Beirne's technical report on vaults~\cite{OBeirne23}
informed the design of BIP-345, which was subsequently
withdrawn~\cite{OB25}. Kirstein et al.~\cite{KLGZ21} formally
verified a single regenerating vault design in Dafny; our work
differs in comparing five designs under a uniform methodology
rather than deeply verifying one.

\paragraph{Covenant proposals and analysis.}
Each covenant proposal has its own design documentation and
community analysis~\cite{BIP119,BIP443,BIP345,BIP347,BIP348}, but no
peer-reviewed work compares them empirically. Harding~\cite{Har24}
published a back-of-envelope estimate of \opvault{} single-block
recovery-cost saturation (historically framed as ``watchtower
exhaustion'') at ${\sim}3{,}000$ splits per block based on
assumed---not measured---transaction sizes; we confirm $3{,}427$ for
\opvault{} and measure $6{,}172$ for \opccv{}, revealing \opccv{}'s
roughly $2\times$ worse unbatched exposure. Poelstra~\cite{Poe21}
demonstrated \opcat{}-based introspection with dynamic destination
and recursive covenants; our \catcsfs{} vault implements a
simplified variant (fixed destination, no recursion) that isolates
the dual-verification security properties. Rijndael~\cite{Rij24}
built a Schnorr-G-trick vault without \opcsfs{}, demonstrating that
\opcat{} alone enables covenant enforcement.

\paragraph{Formal models of Bitcoin contracts.}
The Bitcoin smart-contract formalisation lineage---Atzei et al.'s
formal transaction model~\cite{ABLZ18}, BitML and its
extensions~\cite{BZ18,ABL19,BLM22}, BitML's covenant
extension~\cite{BLZ20}, and the Ethereum-side SoK
methodology~\cite{ABC17}---provides the process-algebraic substrate
against which our Alloy state-machine models stand as a
complementary abstraction level, operating over the $(f, a, g, b)$
design tuple. dgpv~\cite{dgpv24} applied Alloy to Rijndael's
\opcat{} vault; we extend that methodology to all four BIP-specified
designs plus Simplicity under a shared base model, adding a
threat-model layer and a cross-covenant module.
O'Connor~\cite{OC17} provided Coq-verified denotational semantics
for Simplicity; our Simplicity model operates at the state-machine
level, complementing that work with application-level verification.

\paragraph{Layer-two security and rational adversaries.}
Gudgeon et al.'s SoK~\cite{GMRMG20} systematises layer-two
protocols; watchtower-based delegated monitoring has been
formalised in TEE~\cite{LGSM19}, on-chain
arbitration~\cite{McCorry19}, and UC-style state-channel
constructions~\cite{DFH18,DEFM19,ABEFHMM21,Aumayr24}. Our class \classscale{}
analysis shares the always-online monitoring assumption but reframes
the question from ``how does an attacker exhaust the watchtower?''
to ``what is the structural per-UTXO recovery-cost tail of partial
withdrawal, and how does defender batching interact with it?'' On the
rational-adversary side, miner-bribery analyses against timelocked
Bitcoin primitives~\cite{TYME21,WSZN23,NKW21} formalise the
adversary class that covenant vaults defend against during the
unvault delay. Riard~\cite{Ria23} demonstrated replacement cycling
against Lightning HTLCs, a mempool-manipulation class that may
apply to covenant recovery transactions and is out of scope here.

\paragraph{Positioning.}
The qualitative covenants.info matrix~\cite{covenants_info}
compares proposals across feature dimensions without transaction
sizes, attack costs, or fee analysis. No prior work combines
empirical measurement, formal verification, and
fee-environment-dependent security analysis across multiple
covenant vault designs. Our contribution is the integration of
these three threads under a uniform adapter-based methodology
that enables both Proposition~\ref{thm:griefing_safety} and the
design-space framework of Section~\ref{sec:imposs:design_space}.

\providecommand{\todo}[2][]{\textbf{[TODO: #2]}}
